% Part: proof-theory
% Chapter: sequent-calculus
% Section: interpretation-rules

\documentclass[../../../include/open-logic-section]{subfiles}

\begin{document}

\olfileid[pt]{pt}{seq}{irl}

\olsection{Interpretação das regras}

Recorde a interpretação de um sequente: um sequente $\Gamma \Sequent
\Delta$ é verdadeiro em !!a{structure}~$\Struct{M}$ se pelo menos um
$!A \in \Gamma$ for falso ou pelo menos um~$!A \in \Delta$ for
verdadeiro. Podemos ler as regras lógicas de~\Log{G3c} como uma
codificação das condições de verdade de suas !!{formula}s principais.
Considere \RightR{\lif}. Sua !!{formula} principal é $!A \lif !B$. Ela
é verdadeira se $!A$ for falso ou~$!B$ for verdadeiro. Portanto, o
sequente $\ \Sequent !A \lif !B$ é verdadeiro sse $!A \Sequent !B$ for
verdadeiro. Se acrescentarmos as !!{formula}s de contexto $\Gamma$,
$\Delta$ aos antecedentes e sucedentes desses sequentes, obteremos
exatamente a conclusão e a premissa da regra \RightR{\lif}. A situação
é semelhante para~\LeftR{\lif}. $!A \lif !B$ é falsa se $!A$~for
verdadeira e $!B$~for falsa. Assim, $!A \lif !B \Sequent \ $ é
verdadeiro sse tanto $\ \Sequent !A$ quanto $!B \Sequent \ $ forem
verdadeiros. A conclusão de \LeftR{\lif} é equivalente à conjunção de
suas premissas. Esse padrão garante que as regras de~\Log{G3c} sejam
corretas (se todas as premissas forem verdadeiras, a conclusão será
verdadeira). Elas também garantem a completude.

De fato, as condições de verdade de qualquer conectivo verofuncional
podem ser descritas dessa maneira como um conjunto de sequentes. Isso
decorre do fato de que toda função de verdade da lógica clássica pode
ser descrita por !!a{formula} na forma normal conjuntiva: uma conjunção
de disjunções de !!{formula}s atômicas e atômicas negadas. Por exemplo,
considere que $!A \oplus !B$ seja verdadeira sse exatamente uma entre
$!A$ e $!B$ for verdadeira, O.s.\@, ``ou exclusivo''. Então $!A
\oplus !B$ é verdadeira sse $(!A \lor !B) \land (\lnot !A \lor \lnot
!B)$ for verdadeira, e será falsa se $(\lnot !A \lor !B) \land (!A
\lor \lnot !B)$ for verdadeira. Portanto, poderíamos introduzir regras
do cálculo de sequentes para $\oplus$ como
\begin{align*}
&
\Axiom$!A, \Gamma \fCenter \Delta, !B$
\Axiom$!B, \Gamma \fCenter \Delta, !A$
\RightLabel{\LeftR{\oplus}}
\BinaryInf$!A \oplus !B, \Gamma \fCenter \Delta$
\DisplayProof
\\
&
\Axiom$\Gamma \fCenter \Delta, !A, !B$
\Axiom$!A, !B, \Gamma \fCenter \Delta$
\RightLabel{\RightR{\oplus}}
\BinaryInf$\Gamma \fCenter \Delta, !A \oplus !B$
\DisplayProof
\end{align*}
e essas regras também seriam corretas e completas. Suas conclusões são
verdadeiras sse todas as premissas forem verdadeiras.

\begin{prob}
Suponha que $!A \downarrow !B$ seja verdadeira sse nem $!A$ nem~$!B$
for verdadeira (``NÃO-OU'', ou ``NOR''). Quais seriam suas regras de \Log{G3c}?
(Encontre dois sequentes que sejam verdadeiros ao mesmo tempo sse $!A
\downarrow !B$ for verdadeira. Isso fornece a regra
\RightR{\downarrow}. Em seguida, encontre um sequente que seja
verdadeiro sse $!A \downarrow !B$ for falsa. Isso fornece a regra
\LeftR{\downarrow}.)
\end{prob}

Em \Log{G3c}, cada conectivo e quantificador tem exatamente duas regras:
uma regra à esquerda e uma regra à direita. Em \Log{G1c}, porém, há
duas regras \LeftR{\land} e duas regras \RightR{\lor}. A razão é que o
sistema original de Gentzen~\Log{LK}, no qual \Log{G1c} é baseado,
tinha uma variante~\Log{LJ} para uma importante lógica não clássica, a
saber, a lógica intuicionista. Para obter um sistema correto e completo
para a lógica intuicionista, restringe-se o sucedente de qualquer
sequente em~\Log{LJ} a no máximo uma !!{formula}. Isso impossibilita o
uso da regra \RightR{\lor} de~\Log{G3c}, pois sua premissa contém tanto
$!A$ quanto~$!B$ no sucedente. Gentzen, portanto, usou um par de regras
\RightR{\lor} para \Log{LJ} e usou as mesmas regras para~\Log{LK}. De
modo semelhante, a versão intuicionista~\Log{G1i} é simplesmente
\Log{G1c} com todos os sucedentes restritos a no máximo uma
!!{formula}. (\Log{G3c} também tem uma versão intuicionista, mas algumas
de suas regras diferem das regras correspondentes em~\Log{G3c}. Por
exemplo, ela também usa um par de regras \RightR{\lor}.)

É fácil ver que as regras estruturais de~\Log{G1c} são corretas. Por
exemplo, se $\Gamma \Sequent \Delta$ contiver uma !!{formula} falsa à
direita ou uma !!{formula} verdadeira à esquerda, $\Gamma \Sequent
\Delta, !A$ também conterá: a mesma fórmula, seja $!A$ verdadeira ou
falsa, é irrelevante. Assim, se a premissa de~\RightR{\Weakening} for
verdadeira, sua conclusão também será. Observe que, aqui, a recíproca
não é verdadeira. A conclusão pode ser verdadeira porque $!A$ é
verdadeira, e então não se garante que a premissa seja verdadeira.

Por que ter regras estruturais? A contração (\RightR{\Contraction},
\LeftR{\Contraction}) é necessária, por exemplo, para !!{prove} $\ \Sequent
!A \lor \lnot !A$. Se você começar com $!A \Sequent !A$,
primeiro obterá $\Sequent !A, \lnot !A$ por~\RightR{\lnot}. Depois,
poderá usar sucessivamente as duas versões de \RightR{\lor}, uma vez
para obter $!A \lor \lnot !A$ a partir de $\lnot !A$ e outra vez para
obter $!A \lor \lnot !A$ a partir de~$!A$. O resultado é $\ \Sequent
!A \lor \lnot !A, !A \lor \lnot !A$. Assim, para obter $\Sequent !A
\lor \lnot !A$ em~\Log{G1c}, é preciso poder ``contrair'' duas
ocorrências da mesma !!{formula} em !!a{proof}:
\[
\Axiom$!A \fCenter !A$
\RightLabel{\RightR{\lnot}}
\UnaryInf$\fCenter !A, \lnot !A$
\RightLabel{\RightR{\lor}}
\UnaryInf$\fCenter !A, !A \lor \lnot !A$
\RightLabel{\RightR{\lor}}
\UnaryInf$\fCenter !A \lor \lnot !A, !A \lor \lnot !A$
\RightLabel{\RightR{\Contraction}}
\UnaryInf$\fCenter !A \lor \lnot !A$
\DisplayProof
\]
De fato, há sequentes válidos que não poderiam ser !!{prove} em
\Log{G1c} sem enfraquecimento ou contração. Por exemplo, !!a{proof} do
sequente $!A \Sequent !B \lif !A$ em~\Log{G1c} precisa terminar em
\RightR{\lif} (pois apenas $!B \lif !A$ pode ser principal), e isso
exige $!B, !A \Sequent !A$ como premissa. Por sua vez, isso exige
\LeftR{\Weakening} para obtê-la do axioma $!A \Sequent !A$:
\[
\Axiom$!A \fCenter !A$
\RightLabel{\LeftR{\Weakening}}
\UnaryInf$!B, !A \fCenter !A$
\RightLabel{\RightR{\lif}}
\UnaryInf$!A \fCenter !B \lif !A$
\DisplayProof
\]

Em \Log{G3c}, as regras de contração e enfraquecimento não são
necessárias. O enfraquecimento está incorporado aos axiomas. A contração
pode, em geral, ser evitada fundindo em uma só as duas cópias do
contexto $\Gamma$, $\Delta$ nas regras com duas premissas. Tanto
\Log{G1c} quanto \Log{G3c} têm regras desse tipo, chamadas de regras de
``compartilhamento de contexto''. Nas regras com uma premissa, a
contração só é necessária para \RightR{\lor}, \LeftR{\land},
\RightR{\lexists} e~\LeftR{\lforall}. As versões dessas regras
em~\Log{G3c} tornam a contração inteiramente supérflua. Há um sistema
\Log{G2c} (veja \olref{tab:G2c}) no qual as regras com duas premissas
têm contextos independentes. Em \Log{G2c}, a contração é ainda mais
importante do que em~\Log{G1c}.

\subfile{rules-G2c}

\end{document}
